\section*{Bloque B. Aplicaciones cuantitativas y optimización}

\noindent\textit{Preguntas 31 a 65. Evalúan la capacidad matemática para calcular impuestos y subsidios pigouvianos, pérdidas de bienestar, equilibrios en permisos de emisión, provisión óptima de bienes públicos, modelos numéricos de Akerlof y Spence y contratos de incentivos bajo riesgo moral.}

\q{La demanda de mercado de un bien químico es $P = 120 - Q$. El costo marginal privado de producción de la industria competitiva es $CMgP = 20 + Q$. El proceso genera efluentes contaminantes con un costo marginal externo de $CMgE = Q$. ¿Cuál es la cantidad producida en el equilibrio de mercado competitivo sin regulación ($Q_C$)?}
\begin{opts}
    \item $Q_C = 25.0$
    \item $Q_C = 33.3$
    \item $Q_C = 40.0$
    \item $Q_C = 50.0$
\end{opts}

\q{Para el mismo mercado de la pregunta anterior ($P = 120 - Q$, $CMgP = 20 + Q$, $CMgE = Q$), ¿cuál es el nivel de producción socialmente eficiente ($Q^*$)?}
\begin{opts}
    \item $Q^* = 25.00$
    \item $Q^* = 40.00$
    \item $Q^* = 50.00$
    \item $Q^* = 33.33$
\end{opts}

\q{Siguiendo el caso anterior ($CMgE = Q$, $Q^* = 100/3$), ¿cuál debe ser el valor del impuesto pigouviano unitario óptimo ($t^*$) para corregir la falla de mercado?}
\begin{opts}
    \item $t^* = 33.33$
    \item $t^* = 20.0$
    \item $t^* = 50.0$
    \item $t^* = 16.67$
\end{opts}

\q{Para los datos anteriores ($P = 120 - Q$, $CMgP = 20 + Q$, $CMgE = Q$, con $Q_C = 50$ y $Q^* = 33.33$), la Pérdida Irrecuperable de Eficiencia ($DWL$) del mercado no regulado es:}
\begin{opts}
    \item $DWL = 250.00$
    \item $DWL = 833.33$
    \item $DWL = 138.89$
    \item $DWL = 416.67$
\end{opts}

\q{La demanda privada de filtros purificadores de agua es $P = 100 - Q$, con costo marginal constante $c = 20$. Cada filtro genera un beneficio externo a la salud pública de $BMgE = 10$ soles constantes. ¿Qué subsidio pigouviano unitario ($s^*$) debe otorgar el gobierno y cuál será la cantidad óptima consumida?}
\begin{opts}
    \item $s^* = 5$ y $Q^* = 85$
    \item $s^* = 10$ y $Q^* = 90$
    \item $s^* = 10$ y $Q^* = 80$
    \item $s^* = 20$ y $Q^* = 100$
\end{opts}

\q{Una fábrica textil genera efluentes en un canal. Su función de beneficios es $\pi(x) = 120 x - x^2$, donde $x$ es el nivel de efluentes vertidos. Aguas abajo, una piscigranja sufre un daño monetario evaluado en $D(x) = 20 x + 0.5 x^2$. Según el Teorema de Coase, ¿cuál es el nivel de vertimiento eficiente ($x^*$) que acordarán mediante negociación?}
\begin{opts}
    \item $x^* = 20$
    \item $x^* = 60$
    \item $x^* = 40$
    \item $x^* = 30$
\end{opts}

\q{En el caso de Coase anterior, si la fábrica no enfrenta restricciones legales de emisión, vertiría sin negociar hasta $d\pi/dx = 120 - 2x = 0 \implies x_0 = 60$, con beneficio $\pi(60) = 3600$ y daño $D(60) = 20(60) + 0.25(60)^2 = 1200 + 900 = 2100$ (bienestar neto $1500$). En el óptimo $x^* = 40$, $\pi(40) = 3200$ y $D(40) = 800 + 400 = 1200$ (bienestar neto $2000$). La ganancia conjunta generada por negociar es:}
\begin{opts}
    \item $800$ soles
    \item $500$ soles
    \item $100$ soles
    \item $300$ soles
\end{opts}

\q{Dos empresas contaminantes tienen costos marginales de reducción de emisiones $CMgR_1 = 4 r_1$ y $CMgR_2 = 2 r_2$, donde $r_i$ son las toneladas reducidas. La autoridad ambiental exige una reducción agregada total de $r_1 + r_2 = 30$ toneladas mediante un mercado de permisos de emisión transferibles. ¿Cuántas toneladas reducirá cada empresa en el equilibrio de costo mínimo?}
\begin{opts}
    \item $r_1 = 20$ y $r_2 = 10$
    \item $r_1 = 8$ y $r_2 = 22$
    \item $r_1 = 10$ y $r_2 = 20$
    \item $r_1 = 15$ y $r_2 = 15$
\end{opts}

\q{Una comunidad de dos habitantes debe decidir la cantidad de lámparas para el alumbrado público ($G$, bien público puro). Las demandas individuales son $P_1 = 60 - G$ y $P_2 = 90 - G$. El costo marginal de provisión de cada lámpara es constante e igual a $CMg = 50$. Según la condición de Samuelson, ¿cuál es el número socialmente óptimo de lámparas ($G^*$)?}
\begin{opts}
    \item $G^* = 40$
    \item $G^* = 30$
    \item $G^* = 50$
    \item $G^* = 25$
\end{opts}

\q{Para el mismo bien público de la pregunta anterior ($P_1 = 60 - G, P_2 = 90 - G, CMg = 50$), si la provisión se deja a la contribución voluntaria privada de los individuos en un juego simultáneo de Nash, la cantidad total provista será:}
\begin{opts}
    \item $G_{Nash} = 0$
    \item $G_{Nash} = 40$
    \item $G_{Nash} = 10$
    \item $G_{Nash} = 50$
\end{opts}

\q{Siguiendo el caso del bien público óptimo ($G^* = 50$), en un Equilibrio de Lindahl donde cada individuo paga una cuota impositiva ($t_i$) igual a su beneficio marginal individual, ¿cuánto paga cada habitante?}
\begin{opts}
    \item $t_1 = 25$ y $t_2 = 25$
    \item $t_1 = 10$ y $t_2 = 40$
    \item $t_1 = 20$ y $t_2 = 30$
    \item $t_1 = 15$ y $t_2 = 35$
\end{opts}

\q{Un lago comunal permite la extracción de peces. La producción total en función del número de pescadores $L$ es $Y(L) = 200 L - 2 L^2$. El precio del pescado es $P = 1$ y el costo de oportunidad diario de cada pescador es $w = 40$ soles. Compare el número de pescadores bajo libre acceso ($L_{libre}$) con el óptimo social ($L^*$):}
\begin{opts}
    \item $L_{libre} = 60$ y $L^* = 30$
    \item $L_{libre} = 40$ y $L^* = 20$
    \item $L_{libre} = 80$ y $L^* = 40$
    \item $L_{libre} = 50$ y $L^* = 25$
\end{opts}

\q{En el lago de la pregunta anterior, ¿cuál es la renta económica máxima que genera el recurso en el óptimo social ($L^* = 40$) y cuánto se disipa bajo libre acceso ($L_{libre} = 80$)?}
\begin{opts}
    \item Renta óptima $= 6400$ soles y bajo libre acceso se disipa al $100\%$ (renta $= 0$)
    \item Renta óptima $= 3200$ soles y bajo libre acceso se disipa al $100\%$ (renta $= 0$)
    \item Renta óptima $= 1600$ soles y bajo libre acceso la renta es $800$ soles
    \item Renta óptima $= 4800$ soles y bajo libre acceso la renta es $1600$ soles
\end{opts}

\q{En un mercado de autos usados a la Akerlof, existen 50 autos buenos que los vendedores valoran en 8000 dólares y los compradores en 10000 dólares; y 50 autos defectuosos ('lemons') que los vendedores valoran en 2000 dólares y los compradores en 4000 dólares. Los compradores no pueden observar la calidad antes de comprar. ¿Qué autos se comerciarán en equilibrio?}
\begin{opts}
    \item Solo se transan los autos buenos a 10000 dólares sin riesgo de calidad.
    \item No se transa ningún vehículo por colapso completo de la demanda media.
    \item Se transan los 100 vehículos conjuntamente a un precio único de 7000 dólares.
    \item Solo se transan los 50 autos defectuosos entre 2000 y 4000 dólares.
\end{opts}

\q{En el modelo de señalización de Spence, los trabajadores de alta productividad generan un valor de $y_A = 100$ y su costo de educación es $c_A(e) = e$. Los de baja productividad generan $y_B = 60$ y su costo es $c_B(e) = 2e$. ¿Cuál es el rango de años de educación ($e^*$) que permite sostener un equilibrio separador donde solo los trabajadores de alta productividad se educan?}
\begin{opts}
    \item $10 \le e^* \le 30$
    \item $20 \le e^* \le 40$
    \item $40 \le e^* \le 60$
    \item $30 \le e^* \le 50$
\end{opts}

\q{Un propietario tiene una riqueza de 100000 soles y enfrenta un riesgo de incendio con pérdida potencial de 36000 soles. Si ejerce un esfuerzo de prevención $e=1$ (que le cuesta 800 soles), la probabilidad de siniestro es $p = 0.05$. Si no ejerce esfuerzo ($e=0$, costo 0), la probabilidad se eleva a $p = 0.20$. Si una compañía de seguros le ofrece cobertura total a una prima actuarialmente justa con esfuerzo ($1800$ soles), ¿qué esfuerzo elegirá el asegurado?}
\begin{opts}
    \item Elegirá $e=1$ por responsabilidad contractual ante siniestros.
    \item Elegirá esfuerzo intermedio de $e=0.5$ para reducir la prima futura.
    \item Elegirá $e=0$ (riesgo moral), pues está cubierto al 100\%.
    \item Renunciará al seguro por superar la prima su pérdida esperada.
\end{opts}

\q{Considere un modelo Principal-Agente donde el resultado de ventas es $y = e + \varepsilon$, con $\varepsilon \sim N(0, \sigma^2)$. El agente tiene costo de esfuerzo $C(e) = 0.5 c e^2$ y coeficiente de aversión absoluta al riesgo $r$. La comisión por desempeño óptima $b^*$ en un contrato lineal $w = a + b y$ es:}
\begin{opts}
    \item $b^* = \frac{1}{1 + r c \sigma^2}$
    \item $b^* = \frac{r c \sigma^2}{1 + r c \sigma^2}$
    \item $b^* = 1$
    \item $b^* = \frac{1}{r c}$
\end{opts}

\q{En una red de alcantarillado, el beneficio privado de conectarse aumenta con el número de usuarios conectados: $P = 100 - Q + 0.5 Q = 100 - 0.5 Q$, mientras que el costo marginal de instalación es $c = 40$. Si no se consideraba la externalidad de red positiva ($0.5 Q$), la demanda percibida era $P = 100 - Q$. ¿En cuánto se incrementa la cantidad óptima al internalizar la red?}
\begin{opts}
    \item Aumenta de 30 a 60 conexiones
    \item Aumenta de 40 a 80 conexiones
    \item Aumenta de 50 a 100 conexiones
    \item Aumenta de 60 a 120 conexiones
\end{opts}

\q{Para tres usuarios de una planta de tratamiento comunal, las demandas por capacidad $G$ son $P_1 = 30 - G$, $P_2 = 40 - G$ y $P_3 = 50 - G$. El costo marginal de provisión es $CMg = 30$. Según la regla de Samuelson, ¿cuál es la capacidad óptima ($G^*$)?}
\begin{opts}
    \item $G^* = 15$
    \item $G^* = 30$
    \item $G^* = 25$
    \item $G^* = 20$
\end{opts}

\q{Una empresa emisora de gases tiene una función de costo marginal de abatimiento $CMgA = 80 - 2 e$, donde $e$ son las emisiones generadas. Si la autoridad ambiental aplica un impuesto de $t = 30$ soles por tonelada emitida, ¿cuál será el nivel de emisiones ($e^*$) que decidirá emitir la empresa?}
\begin{opts}
    \item $e^* = 40$
    \item $e^* = 10$
    \item $e^* = 25$
    \item $e^* = 15$
\end{opts}

\q{En el mecanismo de Clarke-Groves para decidir si se construye un parque hídrico con costo total de 120 soles compartido en 40 soles por cada uno de 3 vecinos. Las valoraciones netas reportadas son: Vecino 1: $+30$, Vecino 2: $+10$, Vecino 3: $-30$. La suma total es $+10 > 0$ (el parque se construye). ¿Qué impuesto de Clarke debe pagar el Vecino 1?}
\begin{opts}
    \item 0 soles
    \item 20 soles
    \item 30 soles
    \item 10 soles
\end{opts}

\q{En el modelo de salarios de eficiencia de Shapiro y Stiglitz, si el salario de reserva y subsidio de desempleo es $\bar{w} = 30$, el costo del esfuerzo es $e = 10$, la tasa de descuento es $r = 0.05$, la tasa de separación exógena es $b = 0.05$ y la probabilidad de ser detectado holgazaneando es $p = 0.20$. El salario mínimo $w^*$ que previene la holgazanería (No-Shirking Condition) es:}
\begin{opts}
    \item $w^* = 50$
    \item $w^* = 40$
    \item $w^* = 60$
    \item $w^* = 45$
\end{opts}

\q{En el modelo de racionamiento de crédito de Stiglitz y Weiss (1981), cuando un banco eleva la tasa de interés más allá de cierto nivel crítico $r^*$:}
\begin{opts}
    \item la tasa de morosidad crediticia se reduce a cero por autoselección positiva.
    \item todos los demandantes solventes reciben el monto completo de crédito solicitado.
    \item el retorno esperado de la cartera bancaria crece de manera estrictamente lineal.
    \item los proyectos seguros salen y quedan los riesgosos, cayendo el retorno.
\end{opts}

\q{Si una empresa que contamina es simultáneamente un monopolio en el mercado del bien final, la teoría de la segunda mejor opción (Second Best) demuestra que el impuesto pigouviano óptimo:}
\begin{opts}
    \item debe exceder el daño externo ($t > CMgE$) para restringir aún más la escala.
    \item debe ser menor que $CMgE$ (o incluso un subsidio) para no agravar la subproducción.
    \item debe fijarse siempre en cero ante cualquier estructura de demanda lineal.
    \item elimina automáticamente toda pérdida irrecuperable atribuible al poder de mercado.
\end{opts}

\q{Considere una función de bienestar social $W(Q) = 80 Q - Q^2 - 10 Q$, donde el último término representa el daño ambiental externo total y los primeros el excedente privado. La derivada de primer orden $dW/dQ = 0$ arroja un óptimo social de:}
\begin{opts}
    \item $Q^* = 35$
    \item $Q^* = 25$
    \item $Q^* = 40$
    \item $Q^* = 30$
\end{opts}

\q{En la teoría de los Bienes Club de James Buchanan (1965), la provisión de una piscina comunal enfrenta una disyuntiva entre:}
\begin{opts}
    \item la no rivalidad pura y la imposibilidad técnica de cobrar membresía.
    \item repartir costos fijos entre más socios y el costo creciente de congestión.
    \item los impuestos ad-valorem locales y las restricciones cuantitativas de acceso.
    \item la fijación de tarifas Ramsey y la tasa de depreciación del activo fijo.
\end{opts}

\q{Dos empresas vecinas evalúan negociar un acuerdo sobre ruidos molestos. La ganancia bruta conjunta de negociar es de 5000 soles. Sin embargo, los costos legales, peritajes y trámites judiciales (costos de transacción) ascienden a 6500 soles. Según el análisis coaseano:}
\begin{opts}
    \item el regulador expropiará ambas plantas para operarlas bajo propiedad pública.
    \item negociarán exitosamente alcanzando la asignación óptima en sentido de Pareto.
    \item el daño marginal externo se convierte automáticamente en un subsidio legal.
    \item no habrá acuerdo pues los costos de transacción exceden la ganancia conjunta.
\end{opts}

\q{Una EPS evalúa subsidiar la instalación de micromedidores domiciliarios. Cada medidor cuesta 120 soles. El beneficio privado del usuario en ahorro de factura es de 70 soles, mientras que la reducción de fugas genera un beneficio social externo de 60 soles para la cuenca. ¿Qué ocurrirá sin subsidio y qué subsidio mínimo ($s$) induce la adopción eficiente?}
\begin{opts}
    \item Sin subsidio no compra ($70 < 120$); basta un subsidio mínimo de $s = 50$ soles.
    \item El proyecto resulta socialmente ineficiente al superar el costo al beneficio privado.
    \item El subsidio estatal requerido debe cubrir exactamente los 120 soles totales.
    \item El hogar instala voluntariamente el medidor aun sin recibir subsidio alguno.
\end{opts}

\q{Un contrato de seguro médico con copago del $20\%$ implica que si un asegurado consume un tratamiento de rehabilitación que cuesta 2000 soles:}
\begin{opts}
    \item el tratamiento es financiado con fondos de cooperación internacional.
    \item la aseguradora paga los 2000 soles y cobra 400 soles de deducible al año siguiente.
    \item el asegurado paga 400 soles y la aseguradora cubre 1600 soles.
    \item el asegurado paga 2000 soles y la aseguradora nada.
\end{opts}

\q{En el modelo de Akerlof con calidad uniforme de autos $q \sim U[0, 100]$, los vendedores valoran el auto en su calidad $q$ y los compradores en $1.2 q$. Si los compradores pagan el valor esperado de los autos en el mercado, ¿existe algún precio de equilibrio estrictamente positivo?}
\begin{opts}
    \item Sí, cualquier precio positivo entre 0 y 100 vacía el mercado.
    \item No; como $1.2(P/2) = 0.6 P < P$, el mercado colapsa a volumen cero.
    \item Sí, se transan todos los vehículos de calidad al precio $P = 100$.
    \item Sí, existe equilibrio agrupador estable al precio $P = 50$.
\end{opts}

\q{En el análisis de Martin Weitzman (1974) sobre 'Precios vs Cantidades' para regulación ambiental bajo incertidumbre en los costos de abatimiento:}
\begin{opts}
    \item los impuestos (precios) son preferibles si la curva de daño marginal es mucho más empinada que el costo.
    \item ambos instrumentos generan exactamente la misma pérdida esperada bajo cualquier pendiente relativa.
    \item las cuotas (cantidades) son preferibles si la curva de daño marginal es más empinada que el costo.
    \item la incertidumbre en costos favorece siempre estándares tecnológicos rígidos sobre instrumentos de mercado.
\end{opts}

\q{La hipótesis de Charles Tiebout (1956) de 'votar con los pies' sostiene que la provisión eficiente de bienes públicos locales entre municipalidades se logra porque:}
\begin{opts}
    \item los gobiernos locales son auditados continuamente bajo metas presupuestales uniformes.
    \item los servicios públicos locales carecen de costos marginales de congestión o mantenimiento.
    \item el gobierno central asigna presupuestos idénticos per cápita a todas las municipalidades.
    \item los residentes migran al distrito cuya canasta fiscal y servicios coincida con su preferencia.
\end{opts}

\q{En la evaluación económica de daños ambientales globales e inversiones en cuencas, el Informe Stern (2006) utilizó una tasa de descuento social muy cercana a cero ($1.4\%$), lo que provocó que el valor presente de los daños futuros del cambio climático resultara:}
\begin{opts}
    \item sumamente bajo, justificando postergar las inversiones de mitigación.
    \item exactamente igual a cero en términos reales descontados intertemporalmente.
    \item muy elevado, justificando acciones inmediatas y drásticas de mitigación.
    \item totalmente insensible frente al horizonte temporal de evaluación elegido.
\end{opts}

\q{En el modelo de regulación de monopolios con selección adversa de Baron y Myerson (1982), el regulador que no conoce el costo marginal de la empresa debe cederle:}
\begin{opts}
    \item un subsidio directo a los costos fijos únicamente ante riesgo de quiebra.
    \item plena libertad tarifaria sin auditorías contables ni revisión de inversiones.
    \item beneficios económicos negativos garantizados bajo cualquier tipo de tecnología.
    \item una renta informacional positiva para incentivarla a revelar costos bajos.
\end{opts}

\q{Cuando los usuarios residenciales de agua potable pagan una tarifa plana fija mensual sin micromedición (costo marginal de consumo percibido igual a cero):}
\begin{opts}
    \item hay riesgo moral de consumo excesivo hasta anular su beneficio marginal ($BMgP = 0$).
    \item el consumo converge al óptimo donde el beneficio marginal iguala al costo de tratamiento.
    \item la demanda residencial de agua potable se vuelve perfectamente inelástica al ingreso.
    \item el suministro urbano domiciliario adquiere las propiedades de un bien público puro.
\end{opts}

